\element{coordonnees.1}{
  \begin{question}{volumeSpherique}\bareme{1}
    Le volume élémentaire s'écrit en coordonnées sphériques
    \begin{multicols}{2}
      \begin{reponses}
        \bonne{$dV=r^2\sin\theta drd\theta d\varphi$}
        \mauvaise{$dV=r^2\sin\varphi drd\theta d\varphi$}
        \mauvaise{$dV=dr^2\sin\varphi d\theta d\varphi$}
        \mauvaise{$dV=dr^2\cos\varphi d\theta d\varphi$}
      \end{reponses}
    \end{multicols}
  \end{question}
}
\element{coordonnees.1}{
  \begin{question}{volumeCylindrique}\bareme{1}
    La surface élémentaire (pour intégrer sur un cylindre) s'écrit en coordonnées cylindriques
    \begin{multicols}{2}
      \begin{reponses}
        \bonne{$\vec{dS}=rd\theta dz \vec{u_r}$}
        \mauvaise{$\vec{dS}=drdz \vec{u_z}$}
        \mauvaise{$\vec{dS}=dr\sin \theta \vec{u_\theta}$}
        \mauvaise{$\vec{dS}=drd\theta dz \vec{u_r}$}
      \end{reponses}
    \end{multicols}
  \end{question}
}
\setgroupmode{coordonnees.1}{cyclic}

\element{coordonnees.2}{
  \begin{question}{aireSphère}\bareme{1}
    L'aire d'un disque de rayon $R$ est
    \begin{multicols}{4}
      \begin{reponses}
        \bonne{$\pi R^2$}
        \mauvaise{$2\pi R^2$}
        \mauvaise{$4\pi R^3$}
        \mauvaise{$\frac{4}{3}\pi R^3$}
      \end{reponses}
    \end{multicols}
  \end{question}
}
\element{coordonnees.2}{
  \begin{question}{aireCylindre}\bareme{1}
    L'aire latérale d'un cylindre de rayon $R$ et de hauteur $h$ est
    \begin{multicols}{4}
      \begin{reponses}
        \bonne{$2\pi Rh$}
        \mauvaise{$\pi R^2 h$}
        \mauvaise{$2\pi R h^2$}
        \mauvaise{$\pi Rh$}
      \end{reponses}
    \end{multicols}
  \end{question}
}
\setgroupmode{coordonnees.2}{cyclic}


\element{coordonnees.3}{
  \begin{question}{volumeCylindre}\bareme{1}
    Le volume d'un cylindre de rayon $R$ et de hauteur $h$ est
    \begin{multicols}{4}
      \begin{reponses}
        \bonne{$\pi R^2 h$}
        \mauvaise{$2\pi R h^2$}
        \mauvaise{$\pi Rh$}
        \mauvaise{$2\pi Rh$}
      \end{reponses}
    \end{multicols}
  \end{question}
}
\element{coordonnees.3}{
  \begin{question}{volumeBoule}\bareme{1}
    Le volume d'une boule de rayon $R$ est
    \begin{multicols}{4}
      \begin{reponses}
        \bonne{$\frac{4}{3}\pi R^3$}
        \mauvaise{$4\pi R^3$}
        \mauvaise{$\frac{4}{3}\pi R^2$}
        \mauvaise{$4\pi R^2$}
      \end{reponses}
    \end{multicols}
  \end{question}
}
\setgroupmode{coordonnees.3}{cyclic}


\element{coordonnees.4}{
  \begin{question}{bornesSphériques}\bareme{1}
    En coordonnées sphériques :
    \begin{multicols}{2}
      \begin{reponses}
        \bonne{$\theta \in [0,\pi[$ et $\phi \in [0,2\pi[$}
        \mauvaise{$\theta \in [0,2\pi[$ et $\phi \in [0,\pi[$}
        \mauvaise{$\theta \in [0,\pi[$ et $\phi \in [0,\frac{\pi}{2}[$}
        \mauvaise{$\theta \in [0,\pi[$ et $\phi \in [0,\pi[$}
      \end{reponses}
    \end{multicols}
  \end{question}
}
\element{coordonnees.4}{
  \begin{questionmult}{directionOMcylindriques}\bareme{haut=2,p=0}
    En coordonnées cylindriques, $\vec{OM}$ a des composantes selon
    \begin{multicols}{4}
      \begin{reponses}
        \bonne{$\vv{e_r}$}
        \bonne{$\vv{e_z}$}
        \mauvaise{$\vv{e_\theta}$}
        \mauvaise{$\vv{e_\phi}$}
        \mauvaise{$\vv{e_x}$}
        \mauvaise{$\vv{e_y}$}
      \end{reponses}
    \end{multicols}
  \end{questionmult}
}
\element{coordonnees.4}{
  \begin{questionmult}{directionOMspheriques}\bareme{haut=2,p=0}
    En coordonnées sphériques, $\vec{OM}$ a des composantes selon
    \begin{multicols}{4}
      \begin{reponses}
        \bonne{$\vv{e_r}$}
        \mauvaise{$\vv{e_z}$}
        \mauvaise{$\vv{e_\theta}$}
        \mauvaise{$\vv{e_\phi}$}
        \mauvaise{$\vv{e_x}$}
        \mauvaise{$\vv{e_y}$}
      \end{reponses}
    \end{multicols}
  \end{questionmult}
}
\setgroupmode{coordonnees.4}{cyclic}